% GENERATED FILE. Edit canonical lessons, then run npm run generate:books.
% canonical-source-hash: fnv1a64:c915f131703490b8
% canonical-lessons: GU-C228-sojo, GU-C228-sujvu, GU-C228-khanjvalvu, GU-C228-camdi, GU-C228-phollo

\chapter{Swelling and itching}
\label{ch:gu-swelling-and-itching}
\hlchaptermodality{228}

\begin{chapteropening}
\textbf{By the end of this chapter:} \emph{I can say that something is swollen or itchy and ask about a blister.}

Complete the last lesson of chapter 228: I can say that something is swollen or itchy and ask about a blister.
\end{chapteropening}

\section[\texorpdfstring{\gu{સોજો} (sojo)}{સોજો (sojo)}]{\gu{સોજો} (sojo) — swelling}

\label{lesson:GU-C228-sojo}



\begin{quote}
Before the new one: say the Gujarati for to sprain, then the Gujarati for to limp.
\end{quote}

\subsection*{You'll want to know: \gu{સોજો}}
\textbf{\gu{સોજો}} — \emph{sojo} — \textquotedblleft{}swelling\textquotedblright{}.

\begin{quote}
\textbf{\gu{શું} \gu{તમારા} \gu{પગમાં} \gu{સોજો} \gu{છે}?} — \emph{śũ tamārā pagmā̃ sojo che?} — \textquotedblleft{}Is there swelling in your foot?\textquotedblright{}
\end{quote}

A masculine noun.

\begin{cousinweb}
\textbf{\gu{સોજો}} (\emph{sojo}) is a tadbhava, a word inherited from Sanskrit and reshaped by sound change. CDIAL lists it, glossed \textquotedblleft{}a swelling\textquotedblright{}, under \emph{śūyate}, \textquotedblleft{}is swollen\textquotedblright{} (CDIAL 12568).
\end{cousinweb}

\begin{grammarlens}[title={What you've built}]
One more word for everyday life at A2.
\end{grammarlens}

\subsection*{Your turn}
\begin{itemize}
\raggedright
  \item \emph{Say it:} \emph{sojo}
  \item \emph{Say it:} \emph{sojo}, once more
  \item \emph{Say it:} \emph{śũ tamārā pagmā̃ sojo che?} — Is there swelling in your foot?
  \item \emph{Recall:} say the Gujarati for to sprain, then the Gujarati for to limp, then say \emph{sojo} again
\end{itemize}

\begin{tcolorbox}[breakable,colback=teal!4,colframe=teal!35!black,title={Before you move on}]
What does \gu{સોજો} mean? (\textquotedblleft{}Swelling\textquotedblright{}.) Say it once more.
\end{tcolorbox}
\section[\texorpdfstring{\gu{સૂજવું} (sūjvũ)}{સૂજવું (sūjvũ)}]{\gu{સૂજવું} (sūjvũ) — to swell}

\label{lesson:GU-C228-sujvu}



\begin{quote}
Before the new one: say the Gujarati for to limp, then the Gujarati for swelling.
\end{quote}

\subsection*{You'll want to know: \gu{સૂજવું}}
\textbf{\gu{સૂજવું}} — \emph{sūjvũ} — \textquotedblleft{}to swell\textquotedblright{}.

\begin{quote}
\textbf{\gu{મારો} \gu{હાથ} \gu{નથી} \gu{સૂજતો}.} — \emph{māro hāth nathī sūjto} — \textquotedblleft{}My hand is not swelling.\textquotedblright{}
\end{quote}

The swelling itself is \textbf{\gu{સોજો}}.

\begin{cousinweb}
\textbf{\gu{સૂજવું}} (\emph{sūjvũ}) is a tadbhava. CDIAL lists it, glossed \textquotedblleft{}to swell\textquotedblright{}, under Sanskrit \emph{śūyate}, \textquotedblleft{}is swollen\textquotedblright{} (CDIAL 12568).
\end{cousinweb}

\begin{grammarlens}[title={What you've built}]
One more word for everyday life at A2.
\end{grammarlens}

\subsection*{Your turn}
\begin{itemize}
\raggedright
  \item \emph{Say it:} \emph{sūjvũ}
  \item \emph{Say it:} \emph{sūjvũ}, once more
  \item \emph{Say it:} \emph{māro hāth nathī sūjto} — My hand is not swelling.
  \item \emph{Recall:} say the Gujarati for to limp, then the Gujarati for swelling, then say \emph{sūjvũ} again
\end{itemize}

\begin{tcolorbox}[breakable,colback=teal!4,colframe=teal!35!black,title={Before you move on}]
What does \gu{સૂજવું} mean? (\textquotedblleft{}To swell\textquotedblright{}.) Say it once more.
\end{tcolorbox}
\section[\texorpdfstring{\gu{ખંજવાળવું} (khanjvāḷvũ)}{ખંજવાળવું (khanjvāḷvũ)}]{\gu{ખંજવાળવું} (khanjvāḷvũ) — to scratch, to itch}

\label{lesson:GU-C228-khanjvalvu}



\begin{quote}
Before the new one: say the Gujarati for swelling, then the Gujarati for to swell.
\end{quote}

\subsection*{You'll want to know: \gu{ખંજવાળવું}}
\textbf{\gu{ખંજવાળવું}} — \emph{khanjvāḷvũ} — \textquotedblleft{}to scratch, to itch\textquotedblright{}.

\begin{quote}
\textbf{\gu{બાળક} \gu{આંખ} \gu{નથી} \gu{ખંજવાળતું}.} — \emph{bāḷak ā̃kh nathī khanjvāḷtũ} — \textquotedblleft{}The child is not scratching its eye.\textquotedblright{}
\end{quote}

\begin{cousinweb}
\textbf{\gu{ખંજવાળવું}} (\emph{khanjvāḷvũ}) is a tadbhava. CDIAL lists it, glossed \textquotedblleft{}to itch\textquotedblright{}, under the lexical Sanskrit \emph{kharju}, \textquotedblleft{}itching, scratching\textquotedblright{}, a word it takes from Dravidian (CDIAL 3827).
\end{cousinweb}

\begin{grammarlens}[title={What you've built}]
One more word for everyday life at A2.
\end{grammarlens}

\subsection*{Your turn}
\begin{itemize}
\raggedright
  \item \emph{Say it:} \emph{khanjvāḷvũ}
  \item \emph{Say it:} \emph{khanjvāḷvũ}, once more
  \item \emph{Say it:} \emph{bāḷak ā̃kh nathī khanjvāḷtũ} — The child is not scratching its eye.
  \item \emph{Recall:} say the Gujarati for swelling, then the Gujarati for to swell, then say \emph{khanjvāḷvũ} again
\end{itemize}

\begin{tcolorbox}[breakable,colback=teal!4,colframe=teal!35!black,title={Before you move on}]
What does \gu{ખંજવાળવું} mean? (\textquotedblleft{}To scratch, to itch\textquotedblright{}.) Say it once more.
\end{tcolorbox}
\section[\texorpdfstring{\gu{ચામડી} (cāmḍī)}{ચામડી (cāmḍī)}]{\gu{ચામડી} (cāmḍī) — skin}

\label{lesson:GU-C228-camdi}



\begin{quote}
Before the new one: say the Gujarati for to swell, then the Gujarati for to scratch.
\end{quote}

\subsection*{You'll want to know: \gu{ચામડી}}
\textbf{\gu{ચામડી}} — \emph{cāmḍī} — \textquotedblleft{}skin\textquotedblright{}.

\begin{quote}
\textbf{\gu{શું} \gu{તમારી} \gu{ચામડી} \gu{સૂકી} \gu{છે}?} — \emph{śũ tamārī cāmḍī sūkī che?} — \textquotedblleft{}Is your skin dry?\textquotedblright{}
\end{quote}

\begin{cousinweb}
\textbf{\gu{ચામડી}} (\emph{cāmḍī}) is a tadbhava, a word inherited from Sanskrit and reshaped by sound change. It comes from \emph{carman}, \textquotedblleft{}hide, skin\textquotedblright{} (CDIAL 4701).
\end{cousinweb}

\begin{grammarlens}[title={What you've built}]
One more word for everyday life at A2.
\end{grammarlens}

\subsection*{Your turn}
\begin{itemize}
\raggedright
  \item \emph{Say it:} \emph{cāmḍī}
  \item \emph{Say it:} \emph{cāmḍī}, once more
  \item \emph{Say it:} \emph{śũ tamārī cāmḍī sūkī che?} — Is your skin dry?
  \item \emph{Recall:} say the Gujarati for to swell, then the Gujarati for to scratch, then say \emph{cāmḍī} again
\end{itemize}

\begin{tcolorbox}[breakable,colback=teal!4,colframe=teal!35!black,title={Before you move on}]
What does \gu{ચામડી} mean? (\textquotedblleft{}Skin\textquotedblright{}.) Say it once more.
\end{tcolorbox}
\section[\texorpdfstring{\gu{ફોલ્લો} (phollo)}{ફોલ્લો (phollo)}]{\gu{ફોલ્લો} (phollo) — blister, boil}

\label{lesson:GU-C228-phollo}



\begin{quote}
Before the new one: say the Gujarati for to scratch, then the Gujarati for skin.
\end{quote}

\subsection*{You'll want to know: \gu{ફોલ્લો}}
\textbf{\gu{ફોલ્લો}} — \emph{phollo} — \textquotedblleft{}blister, boil\textquotedblright{}.

\begin{quote}
\textbf{\gu{આ} \gu{ફોલ્લો} \gu{મોટો} \gu{નથી}.} — \emph{ā phollo moṭo nathī} — \textquotedblleft{}This blister is not big.\textquotedblright{}
\end{quote}

\begin{cousinweb}
\textbf{\gu{ફોલ્લો}} (\emph{phollo}) is a tadbhava, a word inherited from Sanskrit and reshaped by sound change. It comes from \emph{sphoṭa}, \textquotedblleft{}swelling, boil\textquotedblright{} (CDIAL 13854).
\end{cousinweb}

\begin{grammarlens}[title={What you've built}]
One more word for everyday life at A2.
\end{grammarlens}

\subsection*{Your turn}
\begin{itemize}
\raggedright
  \item \emph{Say it:} \emph{phollo}
  \item \emph{Say it:} \emph{phollo}, once more
  \item \emph{Say it:} \emph{ā phollo moṭo nathī} — This blister is not big.
  \item \emph{Recall:} say the Gujarati for to scratch, then the Gujarati for skin, then say \emph{phollo} again
\end{itemize}

\begin{tcolorbox}[breakable,colback=teal!4,colframe=teal!35!black,title={Before you move on}]
What does \gu{ફોલ્લો} mean? (\textquotedblleft{}Blister, boil\textquotedblright{}.) Say it once more.
\end{tcolorbox}
